% !TEX root = main.tex

%\section{Foundations of Padding-Based Encryption} \label{sec-pgg}

   
\iffalse
The recent NIST post-quantum cryptography standardization contest has  selected a scheme for standardization $\mathsf{CRYSTALS}\mbox{-}\mathsf{Kyber}$ and advanced some other candidates to the fourth round. For focus on the encryption schemes. In the encryption domain, most of the submissions, in particular $\mathsf{CRYSTALS}\mbox{-}\mathsf{Kyber}$, are built by first constructing a \emph{weak} cryptosystem or key-encapsulation mechanism and then applying the Fujisaki-Okamoto (FO) transform~\cite{} to it.  This transform upgrades the weak security to IND-CCA.  FO is essentially a clever form of hybrid encryption in the RO model. 

Unfortunately, an overlooked issue is that \emph{Fujisaki-Okamoto introduces significant ciphertext expansion}. Specifically, to encrypt am $n$-bit message the ciphertext will be about $k+n$ bits, where $k$ is the ciphertext length of the PKE scheme, about $k = 6000$ bits in the case of $\mathsf{CRYSTALS}\mbox{-}\mathsf{Kyber}$. We ask: \emph{Can one have  a better transform in this regard?}
\fi



%Interestingly, there was a similar situation in the 1990s wherein the so-called BR93 IND-CCA scheme~\cite{CCS:BelRog93} had large ciphertext overhead as well.  This overhead was later dropped for RSA-based encryption by  RSA-OAEP~\cite{EC:BelRog94}. Here, ``OAEP'' is an underlying transform that seems quite general.  However, OAEP has not played a role in PQC, perhaps because it is viewed as only applicable only to RSA.  Here, we seek to radically expand its role in PQC.\footnote{Actually, a weakness of OAEP was exposed by Shoup~\cite{C:Shoup01}.  Therefore, we work off of his OAEP+ construct, which is just as efficient.}


%In order to make the overall problem more tractable, our first goal will be to find a suitable abstraction needed in the in the IND-CCA setting.  

\iffalse
\paragraph*{Padding Scheme.}
To explain our approach, we define the notion of a \emph{padding scheme} following~\cite{EC:BelRog94,EC:KilPie09}.
For $\nu, \rs,\ms \in \N$, the associated \emph{padding scheme} is a triple of algorithms $\PADScheme= (\PAD,\PADScheme, \PADScheme^{-1})$ defined as follows.
Algorithm $\PAD$ on input $1^\secp$
outputs a pair $(\pi, \hat{\pi})$ where $\pi: \bits^{\ms+\rs} \to \bits^{\nu}$ and $\hat{\pi}: \bits^{\nu} \to \bits^{\ms} \cup \{\bot\}$ such that $\pi$ is injective and for all $m \in \bits^{\ms}$ and $r \in \bits^\rs$ we have $\hat{\pi}(\pi(m\|r))=m$. 
Algorithm $\PADScheme$ on inputs  $\pi$ and $m \in \bits^{\ms}$ 
outputs $y \in \bits^{\nu}$. Algorithm $\PADScheme^{-1}$ on inputs a mapping $\hat{\pi}$ and $y \in \bits^{\nu}$ 
outputs $m \in \bits^{\ms}$ or $\bot$.
Let $\PADScheme$ be a padding transform from domain $\bits^{\ms + \rho}$ to range $\bits^\nu$.

\iffalse
\paragraph*{Padding-based Encryption with TDFs.}
Let $\calF$ be a TDF with domain $\bits^\nu$.
The associated \emph{padding-based encryption scheme} is a triple of
algorithms $\PADScheme[\calF]=(\kg,\enc,\decc)$ defined in~\figref{fig:pbe}.
\fi

\begin{figure}[t]
\threeCols{0.25}{0.2}{0.2}
    {\underline{$\kg(1^k)$}\\
    $(\pi, \hat{\pi}) \getsr \PAD$ \\
    $(F,F^{-1}) \getsr \TDFIKg(1^k)$ \\
    $\pk \gets (\pi,F)$ \\
    $\sk \gets (\hat{\pi},F^{-1})$ \\
    Return $(\pk,\sk)$ 
    }{
    \underline{$\enc(\pk,m || r)$} \\
    $(\pi,F) \gets \pk$ \\
    $y \gets \pi(m || r)$\\
    $c \gets F(y)$\\
    Return $c$
    }{
    \underline{$\decc(\sk,c)$} \\
    $(\hat{\pi},F^{-1}) \gets \sk$ \\
    $y \gets F^{-1}(c)$\\
    $m  \gets \hat{\pi}(y)$ \\
    %$m\|t \gets m'$\\
    %If $t \neq 0$ then return $\bot$\\
    Return $m$
    } \vspace{-1ex}
\caption{\textbf{Padding based encryption scheme $\PADScheme[\calF]=(\kg,\enc,\decc)$.}}
\label{fig:pbe}
\end{figure}

\iffalse

\paragraph*{OAEP padding scheme.} We  recall the OAEP padding scheme~\cite{EC:BelRog94}. 
Let message length $\mu$, randomness length $\rho$, and redundancy length $\zeta$ be integer parameters, and $\nu = \mu + \rho + \zeta$.
Let $\calG \colon \calK_G \times \bits^\rs \to \bits^{\ms + \zeta}$ and $\calH \colon \calK_H \times \bits^{\ms + \zeta} \to \bits^\rs$ be function families. 
The associated \emph{OAEP padding scheme} is a triple of algorithms $\OAEP[\calG,\calH] = (\calK_\OAEP, \OAEP,\OAEP^{-1})$ defined as follows.  On input $1^\secp$, 
$\calK_\OAEP$ returns $(K_G,K_H)$ where \smash{$K_G\getsr \calK_G(1^\secp)$} and \smash{$K_H\getsr \calK_H(1^\secp)$}, and $\OAEP, \OAEP^{-1}$ are as defined in \figref{fig:oaepscheme}.


\begin{figure}[t]
\twoCols{0.36}{0.36}
{
\textbf{Algorithm} $\OAEP_{(K_G,K_H)}(m \| r)$ \\
%$r \getsr \bits^\rho$ \\
$s \gets  (m \| 0^\zeta) \xorr G(K_G, r)$ \\
$t \gets r \xorr H(K_H, s)$ \\
$x \gets s \| t$\\
Return $x$
}{
\textbf{Algorithm} $\OAEP^{-1}_{(K_G,K_H)}(x)$ \\
$s \| t \gets x$ ; $r \gets t \xorr H(K_H, s)$ \\
$m' \gets s \xorr G(K_G, r)$ \\
If $m'|_\zeta = 0^\zeta$ then return $m'|^\mu$ \\
Return $\bot$  \vspace{0.5ex}
} \vspace{-1ex}
\caption{\textbf{OAEP padding scheme $\OAEP[\calG,\calH]$.}}
\label{fig:oaepscheme}
\end{figure}


\paragraph*{OAEP encryption scheme.}
%Slightly abusing notation, 
As in~\figref{fig:pbe}, we denote by $\OAEP[\calG,\linebreak[1]\calH,\linebreak[1]\calF]=(\OAEP.\kg, \OAEP.\enc,\allowbreak \OAEP.\decc)$ the OAEP-based encryption scheme $\calF$-OAEP with $n = \nu$.  
We typically think of $\calF$ as RSA, and all our results apply to this case under suitable assumptions. 
}



\paragraph*{The Fujisaki-Okamoto Transform.} transformation \cite{C:FujOka99,JC:FujOka13} is a technique to convert weak public-key
encryption schemes into strong ones which resist chosen-ciphertext attack (i.e., are IND-CCA2 secure). Let $\SE=(\calK,\calE,\calD)$ be a private-key encryption scheme and let $\PKE=(\kg,\enc,\dec)$ be a public-key encryption scheme.   Assume $\calK(1^\secp)$ outputs a key $K \in \bits^{\secp}$ and $\PKE.\Coins \subseteq \PKE.\msgsp$.
Moreover, let $\calH \colon \calK_H \times \HDom \to \HRng$ and $\calG \colon \calK_G  \times \PKE.\Coins \to \bits^k$ be hash function families. 
The FO transform $\FO[\calH, \calG,\PKE,\SE] = (\FOKg, \FOEnc, \FODec)$ is defined in \figref{fig:fo-scheme}.

\ifnum\llncs=1
\begin{figure}[t]
\threeCols{0.3}{0.3}{0.38}
    {\underline{$\FOKg(1^k)$}\\
    $(pk',sk') \getsr \PKE.\kg(1^k)$ \\
    $K_H \getsr \calK_H(1^k)$\\
    $K_G \getsr \calK_G(1^k)$\\
    $\pk \gets (pk',K_H,K_G)$ \\
    $\sk \gets (sk',K_H,K_G)$ \\
    Return $(\pk,\sk)$ 
}{
    \underline{$\FOEnc(\pk,m;r)$} \\
    $(pk', K_H, K_G) \gets pk$\\
    $y \gets H(K_H,r)$\\
    $c_1 \gets \PKE.\enc(pk',r;y)$\\
    $K \gets G(K_G,r)$\\
    $c_2 \getsr \calE^{\sy}_K(m)$\\
    $c \gets (c_1,c_2)$\\
    Return $c$
}{
    \underline{$\FODec(\sk,c)$} \\
    $(sk', K_H, K_G) \gets sk$\\
    $r \gets \PKE.\decc(sk',c_1)$ \\
    If $r=\bot$ then return $\bot$\\
    $c_1' \gets \PKE.\enc(\pk',r;H(K_H,r))$\\
    If $c'_1 \neq c_1$ then return $\bot$\\
    $K \gets G(K_G,r)$\\
    $m  \gets \calD^{\sy}_K(c_2)$ \\
    Return $m$ \vspace{0.5ex}
    }
\caption{\textbf{FO transform $\FO[\calH,\calG,\PKE,\SE] = (\FOKg, \FOEnc, \FODec)$.}}
\label{fig:fo-scheme}
\end{figure}
\else 
\begin{figure}
\threeCols{0.24}{0.23}{0.3}
    {\underline{$\FOKg(1^k)$}\\
    $(pk',sk') \getsr \PKE.\kg(1^k)$ \\
    $K_H \getsr \calK_H(1^k)$\\
    $K_G \getsr \calK_G(1^k)$\\
    $\pk \gets (pk',K_H,K_G)$ \\
    $\sk \gets (sk',K_H,K_G)$ \\
    Return $(\pk,\sk)$ 
}{
    \underline{$\FOEnc(\pk,m;r)$} \\
    $(pk', K_H, K_G) \gets pk$\\
    $y \gets H(K_H,r)$\\
    $c_1 \gets \PKE.\enc(pk',r;y)$\\
    $K \gets G(K_G,r)$\\
    $c_2 \getsr \calE^{\sy}_K(m)$\\
    $c \gets (c_1,c_2)$\\
    Return $c$
}{
    \underline{$\FODec(\sk,c)$} \\
    $(sk', K_H, K_G) \gets sk$\\
    $r \gets \PKE.\decc(sk',c_1)$ \\
    If $r=\bot$ then return $\bot$\\
    $c_1' \gets \PKE.\enc(\pk',r;H(K_H,r))$\\
    If $c'_1 \neq c_1$ then return $\bot$\\
    $K \gets G(K_G,r)$\\
    $m  \gets \calD^{\sy}_K(c_2)$ \\
    Return $m$ \vspace{0.5ex}
    }\vspace{-1ex}
\caption{\textbf{FO transform $\FO[\calH,\calG,\PKE,\SE] = (\FOKg, \FOEnc, \FODec)$.}}
\label{fig:fo-scheme}
\end{figure}

\fi 

\fi
\fi
